\section{A product certificate that computes its own bound}\label{sec:budget}
This separate theorem concerns a prospective alternative source. It is
not included in the measured cost claims for the delivered quartics.

Consider a binary Horn calculus on natural-number labels. The derivable
labels form the least set containing zero and closed under the permitted
rule instances $v\leftarrow(p,q)$. Thus zero is true and may be used as a
dummy premise; any other axiom is represented by a rule
$v\leftarrow(0,0)$. Label one is the target. For $N$ instances set
\[
 E=\prod_{i=0}^{N-1}(1+v_i)(1+p_i)(1+q_i),\qquad M=2E+2.
\]
Let $C_0=1$, $C_N=E$, and use the local recurrences
\begin{align}
 C_{i+1}&=C_i(1+v_i)(1+p_i)(1+q_i),\label{eq:budgetrec}\\
 S_{i+1}(1+Mv_i)&=S_i(1+Mp_i)(1+Mq_i),\label{eq:flowrec}
\end{align}
with $S_0=M+1$ and $S_N=1$.

\begin{theorem}[Product budget]
There exist $N$, permitted rule instances, and positive integer streams
satisfying these equations if and only if label one has a finite derivation.
\end{theorem}
\begin{proof}
Telescoping gives $A(M)=B(M)$ for
\[
 A(T)=\prod_i(1+v_iT),\qquad
 B(T)=(1+T)\prod_i(1+p_iT)(1+q_iT).
\]
The coefficients are nonnegative. Their sums are $A(1)\le E$ and
$B(1)\le2E<M$, so every coefficient is a base-$M$ digit. Uniqueness of
base expansion gives equality in $\Z[T]$. Unique factorization into
the normalized factors $1+aT$, for $a>0$, yields exact multiset balance:
\[
 \{v_i:v_i>0\}
 =\{1\}\mathbin{\uplus}\{p_i:p_i>0\}
          \mathbin{\uplus}\{q_i:q_i>0\}.
\]
If the target were outside the least set of derivable labels, count
occurrences outside that set. Every rule with an underivable conclusion
has at least one underivable premise, so the output count is at most
the premise count. The multiset balance makes it one greater, a contradiction.
This counting argument permits cyclic rule dependencies and repeated
instances; it does not require an acyclic ordering of the supplied list.

Conversely, expand the target backwards along a finite proof tree. First
compute $E$ from its labels, then $M$. Let $S_i$ be the product of $1+Ma$
over the currently demanded nonzero labels. These positive integers start
at $M+1$, end at one, and obey \eqref{eq:flowrec}. The running product
gives \eqref{eq:budgetrec}. There is no circular choice: the labels fix
$E$ and $M$ before the streams are assigned.
\end{proof}

This removes a common-multiple condition or an externally supplied power
bound from that flow interface. It does not itself produce a smaller
fixed-arity Diophantine equation. Both finite streams still require
unbounded sequence coding, with all variables and operations paid.
The latest direct compiler already uses an efficient single power bound;
adding the second stream is not assumed to improve its measured size.

